% EAI-compliant structured abstract (PubMed style, <=1000 chars).
\begin{abstract}
\noindent\textbf{Background:} Group-relative RL for LLMs (PPO, GRPO, DPO) is
typically judged at the algorithm-label level, ignoring the
model--reward--sampler--stack combination behind each observation.
\textbf{Methods:} We audit the stack with two triage diagnostics:
Zero-Variance Fraction (ZVF, the share of prompts where all $G$ samples
receive identical reward) and Gradient Utilization $=1-$ZVF.
\textbf{Results:} A fixed first-five-step rule (ZVF$\geq$80\%, reward$\leq$5\%)
triages collapsed runs in a 22-run set with only 2 positive cases
(power-bounded). Online reward and capability diverge: Qwen3-8B gains only
82.0$\to$83.3\% on held-out GSM8K; the two real cross-stack runs (Tinker,
TRL) differ $\approx$17$\times$ in last-10 reward; verl/OpenRLHF are
dry-run placeholders.
\textbf{Conclusions:} The algorithm label alone is an under-specified
experimental treatment; usable signal must be verified per stack.
Code released.
\end{abstract}
